\steppingstone{Undoing Multiplication}
\label{stone:division}

The previous stepping stone ended with a question.

Growing was easy because one part of the multiplication was already simple. One quantity contained only a single kind of piece, so the construction almost explained itself.

Can every multiplication be made to look like growth?

If the answer is yes, then division may not require a new idea at all. It may simply require finding the right companion.

Suppose we wish to multiply

\[
P(a,b)
\]

by another PhiBase quantity.

We are not interested in the whole product.

Not yet.

We care only about the long part.

From the multiplication rule,

\[
P(a,b)P(x,y)
=
P(ax+by,\;ay+bx+by).
\]

The second coordinate is the long part.

Everything becomes simple if that quantity can be made to disappear.

So the problem is no longer multiplication.

The problem is to choose

\[
P(x,y)
\]

so that

\[
ay+bx+by=0.
\]

That is the bridge between multiplication and division.

\section*{Searching For A Playmate}

Try a few examples.

Experiment.

Nothing prevents us from choosing different values for \(x\) and \(y\) and watching what happens.

Eventually one choice refuses to go away.

Let

\[
x=a+b,
\qquad
y=-b.
\]

Now substitute those values into the long part.

\[
\begin{aligned}
ay+bx+by
 &=a(-b)+b(a+b)+b(-b)\\
 &=-ab+ab+b^{2}-b^{2}\\
 &=0.
\end{aligned}
\]

The long part has vanished.

The product now looks remarkably like the simple growth examples from the previous chapter.

We have found the playmate.

Mathematicians already have a name for it.

They call it the \emph{conjugate}.

The conjugate of

\[
P(a,b)
\]

is

\[
P(a+b,-b).
\]

\section*{What Remains?}

Once the long part has disappeared, only the short part is left.

Using the same substitution,

\[
\begin{aligned}
ax+by
 &=a(a+b)+b(-b)\\
 &=a^{2}+ab-b^{2}.
\end{aligned}
\]

Nothing mysterious remains.

The product has become

\[
P(a,b)\,P(a+b,-b)
=
P(a^{2}+ab-b^{2},0).
\]

The remaining integer,

\[
a^{2}+ab-b^{2},
\]

is called the \emph{norm}.

The conjugate removes the long part.

The norm is what is left behind.

\section*{Division}

Now the original question answers itself.

To divide

\[
P(u,v)
\]

by

\[
P(a,b),
\]

multiply both numerator and denominator by the conjugate of the denominator.

The denominator immediately becomes the norm, an ordinary integer.

The numerator is an ordinary PhiBase multiplication.

Division has become the natural undoing of multiplication.

Nothing has been approximated.

Nothing has escaped from PhiBase.

\section*{Builder's Notebook}

Find the conjugate of

\[
P(2,1).
\]

Multiply the pair together.

Verify that the long part is zero.

Then divide

\[
P(8,9)
\]

by

\[
P(2,1)
\]

and confirm that the result is

\[
P(3,2).
\]

\section*{Looking Ahead}

The conjugate quietly introduces negative PhiBase coefficients.

The arithmetic has not changed.

The landscape has become larger.

The next stepping stone explores the path that has been hiding inside that landscape from the very beginning: the Fibonacci Spine.
